%% ====================================================================
%%  abstract.tex
%% ====================================================================
\chapter*{Abstract}
\addcontentsline{toc}{chapter}{Abstract}

Deep learning models like DLoc achieve sub-meter accuracy in indoor WiFi localization by learning a
mapping from Channel State Information (CSI) to spatial probability heatmaps over a floorplan. However,
the large size of DLoc---16.5~million parameters---prevents its use on edge devices such as IoT gateways,
smartphones, or the access points themselves. In this thesis, we investigate whether DLoc's accuracy can
be achieved with a much lighter model. We compress DLoc by training a family of lightweight models:
MobileNetV2, a minimal TinyCNN, a MobileNetV2 encoder with U-Net skip connections, and a state-space
(Mamba) model. Each model is trained both as a standalone deep learning model and as a ``student'' model
distilled from the DLoc teacher. Each model is evaluated along five dimensions: accuracy, robustness to
signal degradation (access-point dropout, Gaussian noise, signal attenuation, and spatial blur),
cross-session generalization, statistical reproducibility, and computational efficiency. The best model,
a MobileNetV2 encoder with U-Net skip connections, reduces the median error from 0.707\,m (DLoc) to
0.640\,m---a 9.5\% improvement---yet requires only 2.41~million parameters ($6.8\times$ fewer) and
$56\times$ fewer floating-point operations than the teacher. Furthermore, when evaluated under the same
leave-one-session-out protocol as the original DLoc paper, the MobileNetV2-UNet model reaches the same
accuracy as DLoc on one fold of the cross-session validation and remains within 0.14\,m of the DLoc
accuracy on the other folds. Knowledge distillation was mainly beneficial to robustness against Gaussian
noise (improving robustness by a factor of 3.4), but only benefited models with sufficient capacity
relative to the data; smaller models that were distilled actually suffered in robustness with the use of
knowledge distillation. Together, these results indicate that it is possible to achieve accurate
localization using lightweight models, and that both architectural and training-based design elements help
to resolve accuracy and robustness issues.

\vspace{0.8em}
\noindent\textbf{Keywords:} WiFi indoor localization, channel state information, knowledge distillation,
model compression, MobileNetV2, U-Net skip connections, robustness, cross-session generalization, edge
deployment.

\clearpage
